\documentclass[11pt]{article}

\usepackage[T1]{fontenc}
\usepackage[utf8]{inputenc}
\usepackage{newtxtext,newtxmath}
\usepackage[letterpaper,margin=0.82in]{geometry}
\usepackage{microtype}
\usepackage{amsmath}
\usepackage{booktabs}
\usepackage{enumitem}
\usepackage{xcolor}
\usepackage{hyperref}
\usepackage{fancyhdr}
\usepackage{titlesec}
\usepackage{tabularx}
\usepackage{longtable}
\usepackage{array}
\usepackage{url}

\definecolor{ink}{HTML}{172033}
\definecolor{accent}{HTML}{315A8A}
\definecolor{soft}{HTML}{EDF3F8}
\definecolor{rulegray}{HTML}{C9D3DE}
\hypersetup{
  colorlinks=true,
  linkcolor=accent,
  citecolor=accent,
  urlcolor=accent,
  pdftitle={CLS Garden Presentations: Exact Scan and Equivalence Closure},
  pdfauthor={Brandon Kirkland},
  pdfsubject={Which CLS census G-matrices conjugate the L-matrices to signed-permutation form, and the equivalence classes of those presentations}
}
\setlength{\parindent}{0pt}
\setlength{\parskip}{0.55em}
\setlist[itemize]{leftmargin=1.35em,itemsep=0.2em,topsep=0.25em}
\setlist[enumerate]{leftmargin=1.5em,itemsep=0.2em,topsep=0.25em}
\titleformat{\section}{\large\bfseries\color{ink}}{}{0pt}{}[\vspace{-0.25em}\color{rulegray}\titlerule]
\titleformat{\subsection}{\normalsize\bfseries\color{accent}}{}{0pt}{}
\pagestyle{fancy}
\fancyhf{}
\lhead{\small CLS Garden presentations: scan and equivalence}
\rhead{\small Pr\'ecis, August 2026}
\cfoot{\small\thepage}
\renewcommand{\headrulewidth}{0.3pt}

\newcommand{\summarybox}[1]{%
  \begin{center}
  \fcolorbox{rulegray}{soft}{\begin{minipage}{0.94\linewidth}\small #1\end{minipage}}
  \end{center}}
\newcommand{\pathcode}[1]{\nolinkurl{#1}}

\begin{document}

\begin{center}
  {\LARGE\bfseries\color{ink} Garden Presentations of the CLS Valise}\par
  \vspace{0.3em}
  {\large Exact scan of the census G-matrix representatives, and closure of the equivalence question}\par
  \vspace{0.55em}
  {\small Pr\'ecis prepared for S.\,J.\ Gates Jr.\ and Y.\ (Tom) Lee}
\end{center}

\summarybox{%
This is the companion to the CLS valise G-matrix census pr\'ecis
(\pathcode{docs/cls-gmatrix-census-precis-for-gates.tex}).  The census found
$123{,}984{,}864$ matrices $G\in\{-1,0,1\}^{12\times 12}$ with $G^2 = A$ and
stored 4{,}536 class representatives (4{,}077 distinct matrices).  The natural
follow-up, prompted by the holoraumy program of \cite{gates2024precis}, is:
\emph{which} stored $G$ conjugate all four CLS color matrices $L_I$ into
signed-permutation (Garden) form?  Exact answer over the stored record:
\textbf{exactly 9 of the 4{,}077}, all first stored in the two
highest-mass census items.  These nine qualifying matrices produce three
Garden-presentation quadruples, and entrywise-verified signed node
relabelings carry every one of them to the CLS presentation with the color
index untouched.  Consequences: exactly \textbf{one signed-node-relabeling
equivalence class} among the stored Garden presentations; their holoraumy
tuples are conjugate, rather than entrywise identical, and carry the same
invariant content.  The nine $G$-matrices are integral, non-monomial basis
changes connecting the CLS monomial presentation to signed node relabelings
of itself.  An internal preliminary extrapolation of $33{,}696$ qualifying
matrices was rejected before delivery because the required taxonomy-constancy
hypothesis is false.}

\section{The question}

The census pr\'ecis enumerated all $G\in\{-1,0,1\}^{12\times 12}$ satisfying
$G^2 = A$ with $A = (L_1+L_2+L_3+L_4)\,L_1^{-1}$ for the CLS color matrices of
\cite[Appendix~C]{gates2024precis}.  The CLS $L_I$ are signed permutations
satisfying the Garden algebra entrywise,
\[
  L_I L_J^{T} + L_J L_I^{T} \;=\; 2\,\delta_{IJ}\,I_{12},
\]
with multiplicative companions $L_I^2 = -I$ for colors 2, 3, 4 (color 1 is the
identity) and anticommutation among colors 2, 3, 4.

Conjugation $C_I = G\,L_I\,G^{-1}$ preserves the multiplicative relations for
any invertible $G$ but destroys monomiality in general.  A \emph{Garden
presentation} here means the ordered quadruple
$C=(C_1,C_2,C_3,C_4)$ when all four $C_I$ are integer signed-permutation
matrices.  A census matrix $G$ \emph{qualifies} when it produces such a
quadruple.  A qualifying $G$ turns the valise back into a monomial adinkra
presentation in a different basis, which is the entry point for comparing
holoraumy data across bases.  The scan asks, exactly, how often that happens
among the stored census representatives and what the resulting presentations
are up to equivalence.

\section{Exact scan of the stored representatives}

Method (\pathcode{scripts/lever_a/garden_presentation_scan.py}, exact
\texttt{Fraction} arithmetic, no floating point anywhere, one core, about
50 seconds):

\begin{itemize}
  \item The four $L_I$ are transcribed in signed-address form from
        \pathcode{src/four_color/cls.rs}
        (Appendix~C of \cite{gates2024precis}), with $L_1 = I_{12}$.
  \item Sanity anchors, all verified entrywise before any scan: the four
        $L_I$ are signed permutations; the Garden algebra above holds;
        $A = (L_1+L_2+L_3+L_4)\,L_1^{-1}$ matches the census artifact
        \pathcode{results/four_color_cls_gmatrix.json} entrywise, tying the
        color matrices to the census target; $L_2^2 = L_3^2 = L_4^2 = -I$
        with $L_1^2 = +I$; colors 2, 3, 4 pairwise anticommute.
  \item Input: all 138 shard files across the three canonical census run
        directories; every stored class representative: 4{,}536 records,
        4{,}077 distinct matrices, 1{,}076 taxonomy keys
        $(\mathrm{nnz}, \mathrm{support}, \mathrm{ranks})$.
  \item Classification by exact rational inverse and matrix product:
        bucket~1 if all four $C_I$ are integer signed permutations, bucket~2
        if all integer but at least one not signed-permutation, bucket~3 if
        any entry is non-integer.
\end{itemize}

\begin{center}
\begin{tabular}{@{}lr@{}}
\toprule
Classification over the 4{,}077 distinct stored matrices & Count \\
\midrule
Qualifying $G$: all four colors signed-permutation & \textbf{9} \\
Integer, only partially signed-permutation & 228 \\
At least one non-integer conjugated color & 3{,}840 \\
\bottomrule
\end{tabular}
\end{center}

Provenance of the nine: 7 are first stored in the shard of census item~52 and
2 in the shard of item~204, the two items contributing the largest census
mass.  All nine also satisfy the Garden algebra entrywise on their conjugated
colors, but this is a corollary, not an extra filter: signed-permutation
conjugates are orthogonal.  For $I=2,3,4$, the similarity-preserved relation
$C_I^2=-I$ therefore gives $C_I^T=C_I^{-1}=-C_I$.  Anticommutation among
colors 2, 3, 4 then gives
$C_I C_J^{T} + C_J C_I^{T}=0$ for distinct $I,J$ in that range.  For a pair
involving color 1, $C_1=I$ and $C_J^T=-C_J$, so
$C_1C_J^T+C_JC_1^T=C_J^T+C_J=0$.  The diagonal relations follow directly
from orthogonality.

\section{Why no population count follows}

Two traps close off any extrapolation from the nine to the full
$123{,}984{,}864$.

\emph{Taxonomy non-constancy.}  The property is not constant on the census
taxonomy classes.  Counts of keys whose stored representatives disagree, by
explicit definition of the classification:

\begin{center}
\small
\begin{tabular}{@{}lr@{}}
\toprule
Definition of the classification & Mixed keys \\
\midrule
Coarse bucket (3-way) & 34 \\
Full four-color signed-permutation flags & 50 \\
Combined (bucket, signed-permutation flags) & 55 \\
Number of integer colors & 182 \\
Full four-color integer flags & 188 \\
\bottomrule
\end{tabular}
\end{center}

Every definition is far above zero, so a taxonomy key does not determine the
conjugation property.  The internal preliminary estimate that $33{,}696$
matrices qualify (extrapolated by assuming constancy on classes) is therefore
unsupported and was rejected before delivery.  Nothing in the shipped census
pr\'ecis ever carried that number; this document records the validated
position.

\emph{Similarity tautologies.}  Multiplicative relations among the $C_I$ are
preserved by any similarity and carry no information beyond
$C_I^2=G L_I^2G^{-1}$.  Thus $C_1^2=+I$ and
$C_2^2=C_3^2=C_4^2=-I$ are forced and are not evidence of an additional
property.  Only the non-similarity-invariant structure (integrality,
monomiality of the conjugated colors, and the transpose algebra, which for
bucket~1 follows from monomiality as shown above) has content.

\section{The nine collapse to three quadruples}

The nine bucket-1 matrices produce only \emph{three} distinct conjugated
quadruples $T^{(1)}, T^{(2)}, T^{(3)}$: seven of the nine land on $T^{(1)}$
and one each on $T^{(2)}$ and $T^{(3)}$.  The collapse mechanism is verified
directly: within each group, $G_j^{-1} G_i$ commutes with all four $L_I$
(commutant collisions), so the different $G$'s differ by elements of the
commutant of the color algebra rather than by genuine presentation changes.

\section{One class: exact signed-monomial equivalence closure}

For ordered presentations $C^{(a)}$ and $C^{(b)}$, the tested adinkra
equivalence relation asks for a signed node map $S$, a color permutation
$\sigma$, and optional color sign rescalings $\epsilon_I\in\{\pm1\}$ such
that
\[
  S C^{(a)}_I S^{-1}=\epsilon_I C^{(b)}_{\sigma(I)}.
\]
Signs carried by the signed monomial matrix $S$ are node signs; they are
distinct from the optional color signs $\epsilon_I$.  The executed result has
$\sigma=\mathrm{id}$ and $\epsilon=(1,1,1,1)$ for every displayed witness:

\begin{center}
\fcolorbox{rulegray}{soft}{\begin{minipage}{0.94\linewidth}\small
All four quadruples (CLS itself, $T^{(1)}$, $T^{(2)}$, $T^{(3)}$) form
\textbf{one equivalence class under signed node relabeling alone}: $\sigma$
is the identity, and no color permutation or color sign rescaling is needed.
The signed node intertwiners are verified entrywise
($S L_I S^{-1}=T_I$ for all four colors).
\end{minipage}}
\end{center}

\noindent In signed-address form (row $i$ carries its single nonzero,
$\mathrm{sign}(a_i)$ in column $|a_i|$):

\begin{center}
\small
\begin{tabular}{@{}ll@{}}
\toprule
Map & Intertwiner $S$ (address form) \\
\midrule
CLS $\to T^{(1)}$ & $(1,\,4,\,2,\,3,\,5,\,-7,\,8,\,-6,\,9,\,-11,\,-12,\,10)$ \\
CLS $\to T^{(2)}$ & $(5,\,-7,\,6,\,-8,\,1,\,4,\,-3,\,-2,\,9,\,-11,\,-12,\,10)$ \\
CLS $\to T^{(3)}$ & $(5,\,-8,\,-7,\,6,\,9,\,11,\,12,\,10,\,1,\,-2,\,4,\,-3)$ \\
\bottomrule
\end{tabular}
\end{center}

The conclusion that the four displayed presentations form one class is
constructive: the three witnesses above already prove it.  The accompanying
search also exhausts the stated signed-monomial equivalence relation.  Each
quadruple is block-diagonal $4{+}4{+}4$, and each 4-node block is a connected
component of the union of its colored node graph.  A signed monomial
intertwiner preserves this colored adjacency and therefore maps each complete
connected block onto a complete connected block.  Its support consequently
contains exactly one $4\times4$ signed-permutation block in each block row and
block column.  The search enumerates all $3!$ block assignments, all
$4!\,2^4=384$ signed permutations inside each matched block, all allowed
global color permutations, and all color sign choices.  Block connectivity,
the resulting plans, and every displayed witness are checked entrywise by
\pathcode{scripts/lever_a/garden_target_equivalence_check.py}.

The block plans make the collapse structural.  Writing each $T$ as a triple
of color permutations of the original CLS blocks:
$T^{(1)}$ uses $(234)/(243)/(243)$, $T^{(2)}$ uses $(23)/(24)/(243)$, and
$T^{(3)}$ uses $(24)/(243)/(34)$.  Within a block family the inducible pure
color permutations are exactly the even ones (the identity and the two
3-cycles); cross-family block conjugators induce exactly the transpositions.
Every plan is realizable by signed node relabeling, which is why everything
collapses to one class.

\section{Holoraumy consequence, and a cross-gadget artifact}

Because the nine qualifying matrices produce presentations related to CLS by
signed node relabelings, their holoraumy tuples are conjugate by the same
relabelings.  They are not generally entrywise identical in the displayed
bases.  There is exactly one holoraumy equivalence class among the stored
Garden presentations, so no new invariant content is present.

One recorded misreading is worth stating precisely, because it is a trap for
any future basis comparison.  The holoraumy gadget
(\pathcode{src/holoraumy.rs}) was evaluated over the four quadruples,
summing traces of products of the $\widetilde{V}$ matrices of one
presentation against those of another.  The resulting cross-matrix,
\[
  \begin{array}{c|cccc}
      & \text{CLS} & T^{(1)} & T^{(2)} & T^{(3)} \\ \hline
    \text{CLS}   & 3 & 0 & 0 & 0 \\
    T^{(1)}      & 0 & 3 & 1 & 1 \\
    T^{(2)}      & 0 & 1 & 3 & 0 \\
    T^{(3)}      & 0 & 1 & 0 & 3
  \end{array}
\]
was briefly read as ``CLS is holoraumy-orthogonal to its conjugates.''  That
reading refutes itself: every target is a conjugate of CLS by construction,
so an isomorphism-invariant comparison must return the self-value 3 on every
entry.  The observed 0s and 1s prove the cross-gadget is a functional of a
pair of \emph{presentations}, not of a pair of isomorphism classes: the trace
of a product of a matrix from one basis against a matrix from another basis
moves under independent conjugation of the two factors.  Direct certificate:
the witness relabeling that carries $T^{(2)}$ onto $T^{(1)}$ moves the
$(T^{(1)}, T^{(2)})$ entry from 1 to the self-value 3.  This basis dependence
applies to every pair of presentations, whether their underlying
representations are equivalent or inequivalent.  The raw cross-gadget is
therefore a basis-alignment functional, not by itself a representation-level
comparison.  Such a comparison requires a common basis convention, an
explicit alignment prescription, optimization over the allowed relabeling
orbit, or a separately proved invariant construction.

What the nine $G$-matrices actually are, then: integral, non-monomial basis
changes connecting the CLS monomial presentation to signed node relabelings
of itself.  The narrow
exact fact that survives is the scan headline: nine of
the 4{,}077 stored representatives conjugate into signed-permutation form,
and that presentation is unique up to signed node relabeling.

\section{What remains open}

Whether \emph{additional} matrices in the nine taxonomy classes (beyond the
stored representatives) also conjugate into Garden form is unknown.  Testing
that requires materializing class members, which needs the census
solution-dump hook and a targeted rerun of the two affected census items,
items~52 and~204, while materializing members of the nine affected taxonomy
classes.  It does not require a new census.  Until that runs, no population
count of any size should be attached to the property.

\section{Artifacts and reproduction}

The computational record is pinned to repository
\url{https://github.com/p1p3dream/adinkra-codespace} at commit
\pathcode{69fe57ba230b0dbd8cbc677b0c1b00b08c5b94b6}.  That commit fixes the
numerical scan, the three target quadruples, the signed-node witnesses, and
the cross-gadget rebasing certificate.  This pr\'ecis and its compiled PDF are
versioned beside the supporting artifacts under \pathcode{docs/}.

\begin{itemize}
  \item \pathcode{results/cls-garden-presentation-scan-20260822.md}:
        full writeup with correction log
  \item \pathcode{scripts/lever_a/garden_presentation_scan.py}:
        the exact scan; \texttt{python3} it, runs in about 50 seconds on one
        core, reproduces $9/228/3{,}840$ and all bookkeeping counts
        (4{,}536 records, 4{,}077 distinct, 1{,}076 keys, 712 multi-representative
        keys, mixed-key counts $34/50/55/182/188$)
  \item Equivalence check, directory \pathcode{scripts/lever_a/}:\\[-0.2em]
        \pathcode{garden_target_equivalence_check.py}.  It performs the
        commutant-collapse checks, connectivity gates, intertwiner
        verification, and complete signed-monomial block search; all checks
        pass
  \item \pathcode{src/four_color/cls.rs}:
        the CLS $L$-matrices in signed-address form
  \item Companion: \pathcode{docs/cls-gmatrix-census-precis-for-gates.tex}
        (the census itself)
  \item This document:
        \pathcode{docs/cls-garden-presentation-equivalence-precis-for-gates.tex}
        and its generated PDF
\end{itemize}

\begin{thebibliography}{99}
\small
\setlength{\itemsep}{0.45\baselineskip}

\bibitem{gates2024precis}
S.\,J.\ Gates Jr.\ and Y.\ Lee,
``A Pr\'ecis: Minimal Four Color Holoraumy and Wolfram's `New Kind
of Science' Paradigm,''
preprint (2024),
\href{https://arxiv.org/abs/2408.09342}{arXiv:2408.09342 [hep-th]}.

\end{thebibliography}

\end{document}
